% !TEX encoding = UTF-8 Unicode
% -*- coding: UTF-8; -*-
% vim: set fenc=utf-8

%%%%%%%%%%%%%%%%%%%%%%%%%%%%%%%%%%%%%%%%%%%%%%%%%%%%%%%%%%%%%%%%%
%% SIMPLE-RESUME-CV
%% <https://github.com/zachscrivena/simple-resume-cv>
%% This is free and unencumbered software released into the
%% public domain; see <http://unlicense.org> for details.
%%%%%%%%%%%%%%%%%%%%%%%%%%%%%%%%%%%%%%%%%%%%%%%%%%%%%%%%%%%%%%%%%

% See "README.md" for instructions on compiling this document.

\documentclass[letterpaper,MMMyyyy,nonstopmode]{simpleresumecv}
% Class options:
% a4paper, letterpaper, nonstopmode, draftmode
% MMMyyyy, ddMMMyyyy, MMMMyyyy, ddMMMMyyyy, yyyyMMdd, yyyyMM, yyyy

%%%%%%%%%%%%%%%%%%%%%%%%%%%%%%%%%%%%%%%%%%%%%%%%%%%%%%%%%%%%%%%%%
%% PREAMBLE.
%%%%%%%%%%%%%%%%%%%%%%%%%%%%%%%%%%%%%%%%%%%%%%%%%%%%%%%%%%%%%%%%%
% CV Info (to be customized).
\newcommand{\CVAuthor}{Yonghun Suh}
\newcommand{\CVTitle}{Yonghun Suh's CV}
\newcommand{\CVNote}{Last updated: %
  \ifcase\the\month
    \or January\or February\or March\or April\or May\or June
    \or July\or August\or September\or October\or November\or December
  \fi \space \the\year}
\newcommand{\CVWebpage}{https://yonghuni.github.io/blog/about/me/}
\newcommand{\CVGithub}{https://github.com/YONGHUNI}
\newcommand{\CVOrcid}{https://orcid.org/0009-0003-6800-2731}


% PDF settings and properties.
\hypersetup{
pdftitle={\CVTitle},
pdfauthor={\CVAuthor},
pdfsubject={\CVWebpage},
pdfcreator={XeLaTeX},
pdfproducer={},
pdfkeywords={},
unicode=true,
bookmarks=true,
bookmarksopen=true,
pdfstartview=FitH,
pdfpagelayout=OneColumn,
pdfpagemode=UseOutlines,
hidelinks,
breaklinks}

% Shorthand.
\newcommand{\Code}[1]{\mbox{\textbf{#1}}}
\newcommand{\CodeCommand}[1]{\mbox{\textbf{\textbackslash{#1}}}}
\makeatletter
\newcommand{\PublicationNumber}[1]{%
  \@ifundefined{pubnum:#1}{??}{\@nameuse{pubnum:#1}}%
}
\newcommand{\PublicationItem}[2]{%
  \phantomsection\label{pub:#1}%
  \expandafter\gdef\csname pubnum:#1\endcsname{#2}%
  \NumberedItem{[#2]}%
}
\newcommand{\PublicationRef}[2]{%
  \hyperref[pub:#1]{Publications, #2 [\PublicationNumber{#1}]}%
}
\makeatother

%%%%%%%%%%%%%%%%%%%%%%%%%%%%%%%%%%%%%%%%%%%%%%%%%%%%%%%%%%%%%%%%%
%% ACTUAL DOCUMENT.
%%%%%%%%%%%%%%%%%%%%%%%%%%%%%%%%%%%%%%%%%%%%%%%%%%%%%%%%%%%%%%%%%

\begin{document}

%%%%%%%%%%%%%%%
% TITLE BLOCK %
%%%%%%%%%%%%%%%

\begin{center}
{\LARGE\bfseries Yonghun Suh}\\[5pt]
Department of Geography, University of Georgia\\
210 Field St., Room 204, Athens, GA 30602, USA\\[4pt]
{\small
\href{mailto:suhyonghun@gmail.com}{\underline{suhyonghun@gmail.com}}
\quad|\quad
\href{\CVOrcid}{\underline{ORCID: 0009-0003-6800-2731}}\\[2pt]
\href{\CVWebpage}{\underline{yonghuni.github.io/blog/about/me/}}
\quad|\quad
\href{\CVGithub}{\underline{github.com/YONGHUNI}}
\quad|\quad
\href{https://scholar.google.com/citations?user=Le1ODbgAAAAJ}{\underline{Google Scholar}}
}
\end{center}

\begin{Body}
\raggedright

%%%%%%%%%%%%%%%
%% EDUCATION %%
%%%%%%%%%%%%%%%

\Section
{Education}
{Education}
{PDF:Education}

\Entry
{\textbf{University of Georgia}},
Athens, GA, USA

\Gap
\textit{       Ph.D. Student in Geography}
\hfill
\DatestampYM{2026}{08} -- Present
\Gap
\SubBulletItem
Advisor: Prof. Weiming Hu

\BigGap
\Entry
{\textbf{Seoul National University}},
Seoul, South Korea

\Gap
\textit{       M.A. in Geography}
\hfill
\DatestampYMD{2021}{09}{02} --
\DatestampYMD{2024}{02}{28}
\Gap
\SubBulletItem
Master's Thesis: ``Real-time Landslide Susceptibility Monitoring Using Spatio-temporal High-resolution Active Remote Sensing Data: An Interpretable Machine Learning Approach''
\Gap
\SubBulletItem
Adviser: Dr. Gunhak Lee
\Gap
\SubBulletItem
GPA: 4.18 / 4.30

\BigGap
\Entry
{\textbf{Kongju National University}},
South Chungcheong Province, South Korea

\Gap
\textit{       B.A. in Geography and B.Sc. in Atmospheric Science}
\hfill
\DatestampYMD{2015}{03}{02} --
\DatestampYMD{2021}{08}{28}

\Gap
\SubBulletItem
GPA: 3.81 / 4.50

%%%%%%%%%%%%%%%
%% RESEARCH INTERESTS %%
%%%%%%%%%%%%%%%

\Section
{Research Interests}
{Research Interests}
{PDF:ResearchInterests}

\Entry
GIScience; GeoAI and spatial machine learning; remote sensing; interpretable modeling;
scalable geospatial computation.
\Gap
\textbf{Application areas}: environmental hazards, weather and climate,
agriculture, and humanitarian spatial analysis.
\BigGap


%%%%%%%%%%%%%%%%%%
%% PUBLICATIONS %%
%%%%%%%%%%%%%%%%%%

\Section
{Publications}
{Publications}
{PDF:Publications}

\SubSection
{Peer-Reviewed Journal Articles}
{Peer-Reviewed Journal Articles}
{PDF:PeerReviewed}

\Gap
\PublicationItem{bike-sharing-air-pollution}{3}
Yoo, E., Roberts, J., \& \textbf{Suh, Y.} (2024).
Delayed effects of air pollution on public bike-sharing system use in Seoul, South Korea: A time series analysis.
\textit{Social Science \& Medicine}, 352, 117030.
\href{https://doi.org/10.1016/j.socscimed.2024.117030}{\underline{doi:10.1016/j.socscimed.2024.117030}}

\BigGap
\PublicationItem{landslide-active-rs}{2}
\textbf{Suh, Y.} \& Lee, G. (2024).
Predicting Landslide Susceptibility Using High-resolution Active Remote Sensing Data:
An Interpretable Machine Learning Approach.
\textit{Journal of the Korean Cartographic Association}, 24(2), 89--111.
[In Korean; English abstract].
\href{https://doi.org/10.16879/jkca.2024.24.2.089}{\underline{doi:10.16879/jkca.2024.24.2.089}}

\BigGap
\PublicationItem{building-scale-population-gwr}{1}
\textbf{Suh, Y.} \& Lee, G. (2023).
Estimation of the de Facto Population at the Building Scale Using a Dasymetric Mapping Method Based on GWR.
\textit{Journal of the Korean Cartographic Association}, 23(1), 21--34.
[In Korean; English abstract].
\href{https://journal.cartography.or.kr/articles/article/Ad7w/}{\underline{Publisher's article page}}

\BigGap
\SubSection
{Preprints}
{Preprints}
{PDF:Preprints}

\Gap
\PublicationItem{shelter-mapping-wash}{1}
Ahn, K.\textsuperscript{\dag}, \textbf{Suh, Y.}\textsuperscript{\dag},
Han, S., Yang, J., Taubenb\"ock, H., \& Cha, M. (2025).
Semi-supervised Shelter Mapping for WASH Accessibility Assessment in Rohingya Refugee Camps.
Preprint (revised March 2026). \href{https://arxiv.org/abs/2511.07231}{\underline{arXiv:2511.07231}}.
\textsuperscript{\dag}Equal first-author contribution.

%%%%%%%%%%%%%%%%%%%%%%%%%
%% RESEARCH EXPERIENCE %%
%%%%%%%%%%%%%%%%%%%%%%%%%

\Section
{Research Experience}
{Research Experience}
{PDF:Research Experience}

\Entry
\textbf{Department of Geography}, University of Georgia, Athens, GA, USA
\Gap
\textit{       Graduate Research Assistant}
\hfill
\DatestampYM{2026}{08} -- Present
\Gap
\SubBulletItem
Contributing computational workflows to research on subseasonal-to-seasonal atmospheric-river forecasting.
\Gap
\SubBulletItem
Developed reproducible Slurm-based HPC workflows for GraphCast Operational weather forecasts (\href{https://github.com/YONGHUNI/s2s-ar-forecasting}{\underline{GitHub: YONGHUNI/s2s-ar-forecasting}}).
\Gap
\SubBulletItem
Faculty advisor: Prof. Weiming Hu.
\BigGap

\Entry
\href{https://www.mpi-sp.org/cha}{\underline{\textbf{Data Science for Humanity}}}, Max Planck Institute for Security and Privacy (MPI-SP), Bochum, Germany
\Gap
\textit{       Research Intern}
\hfill
\DatestampYMD{2025}{05}{01} --
\DatestampYMD{2026}{07}{31}
\Gap
\SubBulletItem
Applied satellite imagery and deep learning to map shelters in Rohingya refugee settlements for WASH accessibility assessment; supported multi-city urban satellite-image processing
\Gap
\SubBulletItem
Co-first-authored an arXiv preprint on semi-supervised shelter mapping and WASH accessibility (\PublicationRef{shelter-mapping-wash}{Preprints})
\Gap
\SubBulletItem
PI: Prof. Meeyoung Cha (MPI-SP \& Korea Advanced Institute of Science and Technology)
\BigGap

\Entry
\href{https://ubairqualitystudy.github.io/EPA-Website/index.html}{\underline{\textbf{UB Clean Air}}}, University at Buffalo, the State University of New York (SUNY), NY, USA
\Gap
\textit{       Project Consultant}
\hfill
\DatestampYMD{2024}{09}{01} --
Present
\Gap
A community-based air quality monitoring project funded by the Environmental Protection Agency
\Gap
\SubBulletItem
Contribute to data collection strategies, spatiotemporal visualization, and modeling workflows
\Gap
\SubBulletItem
Developed interactive dashboards for community-based participatory air monitoring and research dissemination (Conference Presentations [5])
\Gap
\SubBulletItem
PI: Prof. Eunhye Yoo (SUNY)
\BigGap



\Entry
\textbf{Brain Korea Research Team for the Future Landscape}, Seoul National University (SNU), \newline Seoul, South Korea
\Gap
\textit{       Graduate Student Researcher}
\hfill
\DatestampYMD{2021}{09}{01} --
\DatestampYMD{2023}{08}{31}
\Gap
\SubBulletItem
Developed spatial modeling workflows for building-scale population estimation and landslide susceptibility analysis
\Gap
\SubBulletItem
Modeled landslide susceptibility using InSAR displacement, weather radar rainfall estimates, and interpretable machine learning (\PublicationRef{landslide-active-rs}{Peer-Reviewed Journal Articles})
\Gap
\SubBulletItem
Estimated building-scale de facto population in Seoul using dasymetric mapping and geographically weighted regression (\PublicationRef{building-scale-population-gwr}{Peer-Reviewed Journal Articles})
\BigGap


\Entry
\textbf{Development of the Method for Detecting Spatial Interactive Flow Clusters and Its Applicability}, SNU, Seoul, South Korea
\Gap
\textit{       Graduate Student Researcher}
\hfill
\DatestampYMD{2021}{08}{01} --
\DatestampYMD{2022}{06}{30}
\Gap
%https://www.krm.or.kr/krmts/search/detailview/research.html
\SubBulletItem
Reviewed spatial, network, and vector autocorrelation methods for analyzing spatial dependence in interaction data.
\Gap
\SubBulletItem
PI: Prof. Gunhak Lee (SNU)
\BigGap


\Entry
\textbf{Korea Environment Institute (KEI)}, Sejong, South Korea
\Gap
\textit{       Research Intern}
\hfill
\DatestampYMD{2020}{09}{14} --
\DatestampYMD{2020}{12}{14}
\Gap
\SubBulletItem
Processed Sentinel-2 imagery and climate datasets to support research projects on the North Korean environment
\Gap
\SubBulletItem
Assisted Dr. Soojeong Myeong, a Chief Research Fellow, Water and Land Research Group, KEI
\BigGap


%%%%%%%%%%%%%%%%%%
%% CONFERENCE PRESENTATIONS %%
%%%%%%%%%%%%%%%%%%

\Section
{Conference Presentations}
{Conference Presentations}
{PDF:ConferencePresentations}

\NumberedItem{[5]}
\textbf{Suh, Y.} \& Yoo, E. (2025).
From Data to Dissemination: Creating Dashboards for Buffalo African-American Community-based Participatory Air Monitoring.
\textit{Joint Annual Meeting of the International Society of Exposure Science and the International Society for Environmental Epidemiology 2025 (ISES-ISEE 2025)}, Atlanta, USA, Aug. 17–20. \hfill (Poster, \href{https://eventclass.it/isesisee2025/scientific/external-program/session?s=P07#e2609}{Abstract})

\BigGap

\NumberedItem{[4]}
Yoo, E., Lee, J., \& \textbf{Suh, Y.} (2025).
Enhancing Community Air Quality Monitoring: A Spatiotemporal PM2.5 Modeling Approach Using Multiple Data Sources.
\textit{Joint Annual Meeting of the International Society of Exposure Science and the International Society for Environmental Epidemiology 2025 (ISES-ISEE 2025)}, Atlanta, USA, Aug. 17–20. \hfill\href{https://eventclass.it/isesisee2025/scientific/external-program/session?s=TO16#e477}{(Abstract)}

\BigGap

\NumberedItem{[3]}
Yoo, E., Roberts, J., \& \textbf{Suh, Y.} (2024).
Delayed effects of air pollution on public bike-sharing system use in Seoul, South Korea: A time series analysis.
\textit{2024 Spring Conference of Korean Society of Environmental Health}, Gangwon Province, South Korea, May 29–31.

\BigGap

\NumberedItem{[2]}
\textbf{Suh, Y.} \& Lee, G. (2023).
An alternative approach for the landslide prediction using an interpretable machine learning method.
\textit{2023 American Association of Geographers Annual Meeting (2023 AAG)}, Denver, USA, Mar. 23–27.

\BigGap

\NumberedItem{[1]}
\textbf{Suh, Y.} \& Lee, G. (2022).
Estimation of building-scale population density by using a dasymetric-based interpolation method: A case study of Seoul metropolitan area.
\textit{2022 American Association of Geographers Annual Meeting (2022 AAG)}, Virtual, Feb. 25 – Mar. 1.

\BigGap

\SubSection
{In Korean}
{In Korean}
{PDF:PresInKorean}

\BigGap

\NumberedItem{[4]}
\textbf{Suh, Y.} \& Lee, G. (2024).
Real-time Landslide Susceptibility Monitoring Using Spatio-temporal High-resolution Active Remote Sensing Data: An Interpretable Machine Learning Approach.
\textit{2024 Annual Conference of the Korean Geographical Society}, Seoul, South Korea, Jun. 27–28.

\BigGap

\NumberedItem{[3]}
\textbf{Suh, Y.}, Bae, S., \& Jeong, S. (2022).
Accessibility Analysis of Public Cooling Shelters in Seoul Considering Local Temperature: Focusing on the Elderly Population.
\textit{2022 Annual Conference of the Korean Geographical Society}, Seoul, South Korea, Jun. 24–25. \hfill (Poster)

\BigGap

\NumberedItem{[2]}
\textbf{Suh, Y.} \& Lee, G. (2022).
Estimation of the de Facto Population at the Building Scale Using a Dasymetric Mapping Method Based on GWR.
\textit{2022 The Korean Cartographic Association Spring Conference}, Seoul, South Korea, Jun. 11.

\BigGap

\NumberedItem{[1]}
Myeong, S. \& \textbf{Suh, Y.} (2020).
Pre-flood and post-flood damage analysis in the Imjin River Basin.
\textit{2020 Korean Society of Remote Sensing Fall Conference}, Virtual, Nov. 4–6. \hfill (Poster)



%\newpage

%%%%%%%%%%%%%%%%%%%%%%%%%
%% TEACHING EXPERIENCE %%
%%%%%%%%%%%%%%%%%%%%%%%%%

\Section
{Teaching Experience}
{Teaching Experience}
{PDF:Teaching Experience}

\Entry
\textbf{Spatial Analytics 3: Spatio-temporal Data Science}, SNU, Seoul, South Korea
\Gap
\textit{      Teaching Assistant}
\hfill
Spring 2023
\Gap
\SubBulletItem
Provided the entire lab session material using R Markdown and GitHub actions
\Gap
\SubBulletItem
Instructor: Prof. Em. Key-Ho Park (SNU)

\BigGap
\Entry

\textbf{Spatial Analytics 2: Machine Learning}, SNU, Seoul, South Korea
\Gap
\textit{      Teaching Assistant}
\hfill
Fall 2022
\Gap
\SubBulletItem
Conducted TA lab sessions on machine learning algorithms using R
\Gap
\SubBulletItem
Instructor: Prof. Em. Key-Ho Park (SNU)

\BigGap
\Entry

\textbf{Computer Cartography}, SNU, Seoul, South Korea
\Gap
\textit{      Teaching Assistant}
\hfill
Fall 2021
\Gap
\SubBulletItem
Conducted TA lab sessions on cartography \& spatial analytics using ArcGIS and R
\Gap
\SubBulletItem
Instructor: Prof. Gunhak Lee (SNU)

%%%%%%%%%%%%%%%%%%%%%%%%%
%% Work EXPERIENCE %%
%%%%%%%%%%%%%%%%%%%%%%%%%


\Section
{Work Experience}
{Work Experience}
{PDF:Work Experience}

\Entry
\textbf{Office of Information \& Computing Center}, College of Engineering (CoE), SNU, Seoul, South Korea
\Gap
\textit{       System Administrator}
\hfill
\DatestampYMD{2024}{07}{01} --
\DatestampYMD{2025}{03}{31}
\Gap
\SubBulletItem
\textbf{Role}: Managed 2 GPU clusters with multiple nodes (A100 – 4 nodes, GTX 1080 – 100 nodes) for computation demand of 400+ users in CoE, SNU
\Gap
\SubBulletItem
Supported \textit{\href{https://fastmri.snu.ac.kr/}{2024 SNU Fast MRI Challenge}} as a role of MLOps
\BigGap


\Entry
\textbf{The Third Topography Analysis Team}, Third Republic of Korea Army, Gyeonggi Province, South Korea
\Gap
\textit{      Imagery Analysis Specialist}
\hfill
\DatestampYMD{2017}{07}{17} --
\DatestampYMD{2019}{04}{02}
\Gap
\SubBulletItem
\textbf{Role}: Provided crucial terrain information to assist commanders in decision-making
\Gap
\SubBulletItem
Conducted topographical analysis \& provided information through paper and digital maps using ArcGIS and TerraExplorer.
\Gap
\SubBulletItem
Fulfilled mandatory enlisted military service; led the enlisted members of the team as Squad Leader and discharged as Sergeant


%%%%%%%%%%%%%%%%%%%%%%%%%%%
%% AWARDS & SCHOLARSHIPS %%
%%%%%%%%%%%%%%%%%%%%%%%%%%%

\Section
{Awards \&\newline
Scholarships}
{Awards \& Scholarships}
{PDF:AwardsAndScholarships}


\BulletItem
Future Vision Scholarship, SNU
\hfill
\DatestampYM{2023}{02}
\Gap
\SubBulletItem
Received the Scholarship from SNU for outstanding academic performance and exemplary conduct.
\BigGap

\BulletItem
Best Poster Award in Student Poster Competition, Korean Geographical Society (KGS)
\hfill
\DatestampYM{2022}{06}
\Gap
\SubBulletItem
2022 Annual Conference of the KGS
\BigGap

\BulletItem
Outstanding Service Award, SNU
\hfill
\DatestampYM{2022}{02}
\Gap
\SubBulletItem
Awarded by the Brain Korea Research Team for the Future Landscape at SNU for significant contributions to the department.
\BigGap

\BulletItem
Talent Development Scholarship, Jeju International Scholarship Foundation (JISF)
\hfill
\DatestampYM{2016}{05}
\Gap
\SubBulletItem
Awarded by JISF for outstanding academic performance
\BigGap

\BulletItem
The Kongju National University Alumni Association Scholarship
\hfill
\DatestampYM{2015}{11}


%\DatestampY{20(expected)}
%\Gap
%\SubBulletItem


%%%%%%%%%%%%%%%
%% TECHNICAL SKILLS %%
%%%%%%%%%%%%%%%

\Section
{Technical Skills}
{Technical Skills}
{PDF:TechnicalSkills}

\BulletItem
Geospatial \& Remote Sensing
\SubBulletItem
Experienced with: \textbf{ArcGIS} variants, \textbf{QGIS}, \textbf{Google Earth Engine}
\SubBulletItem
Working knowledge: \textbf{ERDAS IMAGINE}, \textbf{HyP3 (ASF)}

\Gap
\BulletItem
Spatial Modeling \& Programming
\SubBulletItem
Advanced: \textbf{R}
\SubBulletItem
Intermediate: \textbf{Bash}, \textbf{Fortran}, \textbf{Python}
\SubBulletItem
Working knowledge: \textbf{Julia}, \textbf{C++}, \textbf{CUDA}

\Gap
\BulletItem
Research Computing \& HPC
\SubBulletItem
Experience with Slurm-managed HPC environments, Linux servers, and GPU clusters
\SubBulletItem
Tools: Slurm, \href{https://www.ansible.com/}{Ansible}, \href{https://kubernetes.io/}{Kubernetes}, \href{https://maas.io/}{Metal as a Service (MaaS)}

\Gap
\BulletItem
Reproducibility \& Publishing
\SubBulletItem
\href{https://git-scm.com/}{Git}, Nix (flakes), \href{https://www.latex-project.org/}{LaTeX}, \href{https://quarto.org/}{Quarto}, R Markdown, GitHub Actions


%%%%%%%%%%%%%%%%%%%%%%%%%
%% OTHER EXPERIENCE %%
%%%%%%%%%%%%%%%%%%%%%%%%%



\Section
{Other Experience}
{Other Experience}
{PDF:Other Experience}

\Entry
\textbf{\href{https://github.com/bacchus-snu}{Bacchus – System Administrator Club}}, \href{https://cse.snu.ac.kr/en/about/student-clubs?selected=Bacchus}{Dept. of Computer Science and Engineering, SNU}, Seoul, South Korea
\Gap
\textit{      Club Member}
\hfill
\DatestampYM{2024}{03} -- \DatestampYM{2025}{04}%Present
\Gap
\SubBulletItem
Set up and maintained Debian package caching servers using Caddy and Reprepro on a Kubernetes cluster
\Gap
\SubBulletItem
Migrated the club's manual webpage from Cloudflare's hosting service to the on-premise server
\BigGap

\Entry
\textbf{2024 Accelerator Programming Winter School}, SNU, Gyeonggi Province, South Korea
\Gap
\textit{      Participant}
\hfill
\DatestampYM{2024}{02}
\Gap
\SubBulletItem
An intensive course covering CUDA programming.
\Gap
\SubBulletItem
Conducted a team project optimizing the inference performance of the GRU model by porting CPU code to GPU kernels.
\BigGap

\Entry
\textbf{Server Management}, Dept. of Geography, SNU, Seoul, South Korea
\Gap
\textit{      System Administrator}
\hfill
\DatestampYM{2022}{11} --
\DatestampYM{2024}{02}
\Gap
\SubBulletItem
Set up a Windows HPC server that includes WSL to help more students utilize the resource.
\Gap


%\newpage
%%%%%%%%%%%%%%%%%%%%%%%%%%%%%%%%%%%%%%%%%%%%
%% SERVICE %%
%%%%%%%%%%%%%%%%%%%%%%%%%%%%%%%%%%%%%%%%%%%%

\Section
{Academic Service}
{Academic Service}
{PDF:AcademicService}

\Entry
\textbf{Journal of the Korean Cartographic Association}, Seoul, South Korea
\Gap
\textit{       Editorial Assistant}
\hfill
\DatestampYMD{2022}{01}{01} --
\DatestampYMD{2022}{12}{31}
\Gap
\SubBulletItem
\textbf{Role}: Managed submissions, reviews, and proofreading processes.
\Gap
\SubBulletItem
Contributed to the editorial activities for three volumes of the journal.
\BigGap

\BigGap

\Entry
\textbf{Graduate Student Council}, Dept. of Geography, SNU, Seoul, South Korea
\Gap
\textit{      Treasurer of the Council}
\hfill
\DatestampYMD{2022}{09}{01} --
\DatestampYMD{2023}{08}{31}
\Gap
\SubBulletItem
Coordinated departmental events and managed a graduate field trip



\end{Body}

%%%%%%%%%%%
% CV NOTE %
%%%%%%%%%%%

\UseNoteFont%
\null\hfill%
[\textit{\CVNote}]

\end{document}
